\section{\rda}
\label{sec:rda}

The core static data-flow analysis is \rda, which is a combination of a static value analysis that tracks all possible values for each variable (in binary analysis, variables refer to value containers, such as registers, memory cells, etc.) at each program location and a data dependency analysis that tracks where each value is defined and used.
This is a critical difference from prior work:
Both \karonte and \satc use taint-tracking, which means they can only mark certain variables and values as \emph{tainted} by input data or not.
In comparison, \rda creates def-use graphs, with which \mango can reason about how any value was created from its source values and what operations have been applied.
We call these expressions \emph{\richexprs} because they capture not only the concrete values but also rich information about their history and dependencies.

An example \richexpr is: \texttt{"df -d " .. itoa(atoi(input\_1)) .. " -h " .. input\_0)} where ``\texttt{..}'' represents the string concatenation operation.
In this case, the underlying \richexpr can be broken into its 4 components, \texttt{"df -d "}, \texttt{itoa(atoi(input\_1))}, \texttt{" -h "}, and \texttt{input\_0}, at any time with their corresponding dependencies already held in a global dependency graph.
% Figure~\ref{fig:richexpr} illustrates the difference between \mango and taint-analysis-based solutions.
If a call to \texttt{system()} has a \richexpr as the first argument (such as the one described), \mango can determine that \texttt{input\_1} does not lead to arbitrary command injection as its value was derived from an integer. 
In contrast, \texttt{input\_0} may lead to a command injection vulnerability as its origin value is unhindered.
Neither \karonte nor \satc could determine that \texttt{input\_1} would not contribute to a vulnerability, which leads to higher false positive rates and reduced analysis speed due to over-tainting.

\rda is the foundation for various sub-analyses on which \mango depends, such as calling convention inference and function prototype inference.
    \mango uses calling convention inference to determine the ABI of each function, which can be different even for functions within the same binary.
    \mango further uses function prototype inference to determine the number and the sizes of arguments for each function and if there is a return value.
    Both types of information are necessary for inter-function analysis.
    These sub-analyses are purely engineering work of existing research~\cite{balakrishnan2004analyzing,zhang2007parameter} and are not our scientific contributions in this paper.

\subsection{\rda Basics}

\noindent
\textbf{Value domain.}
\rda employs a value domain that is specifically designed for finding command injection vulnerabilities.
Basic value types in this domain include bit-vectors (integers and strings with fixed lengths) and a TOP ($\top$) value.
Each value in this domain is either a \emph{MultiValue} that contains at most $N$ bit-vector values, a $\top$ value, or an \emph{Expression} that is composed of \emph{operations} that are applied on MultiValues.
Each instance of MultiValue and Expression is of a fixed number of bits.
\rda is flow-sensitive but not path-sensitive, which means that whenever multiple paths merge, the values in each variable will be merged, creating MultiValues with more than one basic value inside.
A MultiValue with more than $N$ bit-vectors becomes a $\top$, and overly deep expressions (we empirically limit it to 20) also become $\top$.
We choose a small $N$ ($N=5$) to guarantee fast convergence of the data-flow analysis.

\smallskip
\noindent
\textbf{Data dependency tracking.}
The data dependency analysis in \rda is similar to traditional, source-code-based data dependency analysis:
\rda tracks where each variable is first defined and keeps a separate record for program locations where each defined variable is used.
The definition record is associated with each value (bit-vectors or $\top$ values).
During path merging, \rda will adjust the association of definition records and re-associate definition records of old values to newly created $\top$ values if needed.

\smallskip
\noindent
\textbf{Memory model.}
The abstract state used in \rda models ``loads from'' and ``stores into'' for registers, global memory regions, stack regions, and heap locations in a bit-precise manner.
Because \rda performs intra-function analysis by default, it only models one stack frame (also known as activation record).

\subsection{Inter-function Analysis}
To resolve sinks, \mango requires inter-function analysis.
Building an inter-function \rda from an intra-function \rda is straightforward.
At every call site, \mango makes a copy of the abstract analysis state thus far, creates a new stack frame to simulate the call, invokes a new \rda on the copy of the abstract state, and then waits until the new \rda instance terminates.
Then, \mango merges all exit states that the new \rda yields to get a new output state and uses that state to resume the outer-level \rda analysis.

\subsection{Function Summaries}
\label{sec:function_handlers}

The more accurately \rda models the execution of a binary, the higher level of precision \mango can achieve.
% Interestingly, while both \karonte and \satc model multi-binary interactions, they do not model the exact execution artifacts of said binary's libraries, which include libc.
\karonte and \satc are based on taint tracking, which models dependency relationships between input arguments and output values of certain library functions.
Precisely propagating taints through many library functions (especially in \texttt{glibc}) can be difficult due to the complexity of the logic and implementations of these library functions, which may lead to excessive amounts of false positives and false negatives.
Therefore, \mango implements function summaries for common library functions to ensure that values and data dependencies are modeled as accurately as possible.

Re-implementing library functions has the added benefit of reducing false positives caused by changes in value types.
It is a common and normal execution pattern for programs to manipulate strings into non-injectable forms, such as using \texttt{atoi} to transfer the string into an integer or the \texttt{"\%d"} format specifier for the \texttt{printf} function family, making injections impossible.
\satc does not recognize these input type changes and will report a (false-positive) vulnerability even if the tainted value is not injectable.
%Without proper value modeling it is easy to over-taint tracked values.
